\chapter{Connected transversals of 5-colorings}
Kriesell and Mohr \cite{matthias_2022} showed that all graphs on five vertices and at most 6 edges has (*) property. 
We will try to extend this result to 7 edges. First let us see what are those graphs on five vertices and 7 edges.
Let us assume the graphs are connected, otherwise by Theorem 1 in \cite{matthias_2022}. we have that the graph has (*) property.
Now let's see those graphs.

\vspace{0.1cm}
\includegraphics[width=13.5cm]{img/7-edge-cases.eps}

The (c) graph has (*) property since $K_4$ has (*) property (By theorem 4 in \cite{matthias_2022}), and $K_4 + 1$ vertex with degree 1 also has (*) property (By lemma 1 in \cite{matthias_2022}).

We will show that the minimal in terms of |V| + |E| counter example for graph (a) to not hold property (*) must be 3-connected.

\begin{lemma}
    A minimal counter example such that graph (a) doesn't hold property (*) is 2-connected.
\end{lemma}
\begin{proof}
    Let's denote graph (a) as $H$, Let $G$ be the minimal graph such that graph $H$ is isomorphic to $H(G, \mathcal{C}, T)$, where transversal $T := \{1, 2, 3, 4, 5\}$, and $\mathcal{C}$ is the coloring of $G$
    The minimality implies that $(\forall A,B \in \mathcal{C} : A \neq B)$ $G[A \cup B]$ is a single path between the transversal vertices of corresponding colors from $T$.

    Assume for contradiction that $G$ is 1-connected. Then there exists a cut vertex in $G$, let us denote it as $x$. We have six cases to consider.
    
    \vspace{0.3cm}
    \includegraphics[width=13cm]{img/hourglass+1-cases.eps}
    \vspace{0.3cm}


    For case (a) we have that the vertex on the right side of the cut has degree at least 2, but it can have Kempe chain only to the $x$, hence it's a contradiction.
    For case (b) we can contract the edges of the right side of the cut and still get the counter example, hence the graph is not minimal, a contradiction.
    For case (c) the vertex $x$ should have color $1$, otherwise if it is the vertex with color $1$ is on any side of the cut, we wouldn't be able to 
    have Kempe chain to the 2 other transversal vertices of the opposite side of the cut. Now we have that $x$ has color $1$ and WLOG $2$ is on the left side of the cut. 
    And the other vertex in that same side of the cut must be of color $3$ or $4$, it's not $4$ since it has degree 3, and it wouldn't be able to connect to vertex $5$, hence
we have the following graph.

\vspace{0.3cm}
\includegraphics[width=13.5cm]{img/hourglass+1-cases-2-conn(c).eps}
\vspace{0.1cm}


But now we see that $2$ and $4$ can not have a Kempe chain between each other, hence they can not be connected in $H$, hence it's a contradiction.
For the case (d) we have that the vertex on the right side of the cut has degree at least 2, hence the color of $x$ is the color of the vertex itself,
so we can contract the paths from it to $x$ and still get the counter example, hence a contradiction since the graph wasn't minimal.
For the case (e) on the right side of the cut we have that both of the vertices have degree at least 2, hence there should pass at least 3 paths from right 
side of the cut to the left side of the cut, in this case the color of the vertex $x$ should be 1, and any combination of colors of vertices on right side 
still forces to not have a Kempe chain between at least one vertex on the right side and one vertex on the left side, hence it's a contradiction.
For the case (f) we could just contract all the vertices on the right side of the cut to vertex $x$ and still get the counter example, hence it's a contradiction.
Hence the graph $G$ is 2-connected.
\end{proof}
